\IfFileExists{stacks-project.cls}{%
\documentclass{stacks-project}
}{%
\documentclass{amsart}
}

% For dealing with references we use the comment environment
\usepackage{verbatim}
\newenvironment{reference}{\comment}{\endcomment}
%\newenvironment{reference}{}{}
\newenvironment{slogan}{\comment}{\endcomment}
\newenvironment{history}{\comment}{\endcomment}

% For commutative diagrams we use Xy-pic
\usepackage[all]{xy}

% We use 2cell for 2-commutative diagrams.
\xyoption{2cell}
\UseAllTwocells

% We use multicol for the list of chapters between chapters
\usepackage{multicol}

% This is generally recommended for better output
\usepackage{lmodern}
\usepackage[T1]{fontenc}

% For cross-file-references
\usepackage{xr-hyper}

% Package for hypertext links:
\usepackage{hyperref}

% For any local file, say "hello.tex" you want to link to please
% use \externaldocument[hello-]{hello}
\externaldocument[introduction-]{introduction}
\externaldocument[conventions-]{conventions}
\externaldocument[sets-]{sets}
\externaldocument[categories-]{categories}
\externaldocument[topology-]{topology}
\externaldocument[sheaves-]{sheaves}
\externaldocument[sites-]{sites}
\externaldocument[stacks-]{stacks}
\externaldocument[fields-]{fields}
\externaldocument[algebra-]{algebra}
\externaldocument[brauer-]{brauer}
\externaldocument[homology-]{homology}
\externaldocument[derived-]{derived}
\externaldocument[simplicial-]{simplicial}
\externaldocument[more-algebra-]{more-algebra}
\externaldocument[smoothing-]{smoothing}
\externaldocument[modules-]{modules}
\externaldocument[sites-modules-]{sites-modules}
\externaldocument[injectives-]{injectives}
\externaldocument[cohomology-]{cohomology}
\externaldocument[sites-cohomology-]{sites-cohomology}
\externaldocument[dga-]{dga}
\externaldocument[dpa-]{dpa}
\externaldocument[sdga-]{sdga}
\externaldocument[hypercovering-]{hypercovering}
\externaldocument[schemes-]{schemes}
\externaldocument[constructions-]{constructions}
\externaldocument[properties-]{properties}
\externaldocument[morphisms-]{morphisms}
\externaldocument[coherent-]{coherent}
\externaldocument[divisors-]{divisors}
\externaldocument[limits-]{limits}
\externaldocument[varieties-]{varieties}
\externaldocument[topologies-]{topologies}
\externaldocument[descent-]{descent}
\externaldocument[perfect-]{perfect}
\externaldocument[more-morphisms-]{more-morphisms}
\externaldocument[flat-]{flat}
\externaldocument[groupoids-]{groupoids}
\externaldocument[more-groupoids-]{more-groupoids}
\externaldocument[etale-]{etale}
\externaldocument[chow-]{chow}
\externaldocument[intersection-]{intersection}
\externaldocument[pic-]{pic}
\externaldocument[weil-]{weil}
\externaldocument[adequate-]{adequate}
\externaldocument[dualizing-]{dualizing}
\externaldocument[duality-]{duality}
\externaldocument[discriminant-]{discriminant}
\externaldocument[derham-]{derham}
\externaldocument[local-cohomology-]{local-cohomology}
\externaldocument[algebraization-]{algebraization}
\externaldocument[curves-]{curves}
\externaldocument[resolve-]{resolve}
\externaldocument[models-]{models}
\externaldocument[functors-]{functors}
\externaldocument[equiv-]{equiv}
\externaldocument[pione-]{pione}
\externaldocument[etale-cohomology-]{etale-cohomology}
\externaldocument[proetale-]{proetale}
\externaldocument[relative-cycles-]{relative-cycles}
\externaldocument[more-etale-]{more-etale}
\externaldocument[trace-]{trace}
\externaldocument[crystalline-]{crystalline}
\externaldocument[spaces-]{spaces}
\externaldocument[spaces-properties-]{spaces-properties}
\externaldocument[spaces-morphisms-]{spaces-morphisms}
\externaldocument[decent-spaces-]{decent-spaces}
\externaldocument[spaces-cohomology-]{spaces-cohomology}
\externaldocument[spaces-limits-]{spaces-limits}
\externaldocument[spaces-divisors-]{spaces-divisors}
\externaldocument[spaces-over-fields-]{spaces-over-fields}
\externaldocument[spaces-topologies-]{spaces-topologies}
\externaldocument[spaces-descent-]{spaces-descent}
\externaldocument[spaces-perfect-]{spaces-perfect}
\externaldocument[spaces-more-morphisms-]{spaces-more-morphisms}
\externaldocument[spaces-flat-]{spaces-flat}
\externaldocument[spaces-groupoids-]{spaces-groupoids}
\externaldocument[spaces-more-groupoids-]{spaces-more-groupoids}
\externaldocument[bootstrap-]{bootstrap}
\externaldocument[spaces-pushouts-]{spaces-pushouts}
\externaldocument[spaces-chow-]{spaces-chow}
\externaldocument[groupoids-quotients-]{groupoids-quotients}
\externaldocument[spaces-more-cohomology-]{spaces-more-cohomology}
\externaldocument[spaces-simplicial-]{spaces-simplicial}
\externaldocument[spaces-duality-]{spaces-duality}
\externaldocument[formal-spaces-]{formal-spaces}
\externaldocument[restricted-]{restricted}
\externaldocument[spaces-resolve-]{spaces-resolve}
\externaldocument[formal-defos-]{formal-defos}
\externaldocument[defos-]{defos}
\externaldocument[cotangent-]{cotangent}
\externaldocument[examples-defos-]{examples-defos}
\externaldocument[algebraic-]{algebraic}
\externaldocument[examples-stacks-]{examples-stacks}
\externaldocument[stacks-sheaves-]{stacks-sheaves}
\externaldocument[criteria-]{criteria}
\externaldocument[artin-]{artin}
\externaldocument[quot-]{quot}
\externaldocument[stacks-properties-]{stacks-properties}
\externaldocument[stacks-morphisms-]{stacks-morphisms}
\externaldocument[stacks-limits-]{stacks-limits}
\externaldocument[stacks-cohomology-]{stacks-cohomology}
\externaldocument[stacks-perfect-]{stacks-perfect}
\externaldocument[stacks-introduction-]{stacks-introduction}
\externaldocument[stacks-more-morphisms-]{stacks-more-morphisms}
\externaldocument[stacks-geometry-]{stacks-geometry}
\externaldocument[moduli-]{moduli}
\externaldocument[moduli-curves-]{moduli-curves}
\externaldocument[examples-]{examples}
\externaldocument[exercises-]{exercises}
\externaldocument[guide-]{guide}
\externaldocument[desirables-]{desirables}
\externaldocument[coding-]{coding}
\externaldocument[obsolete-]{obsolete}
\externaldocument[fdl-]{fdl}
\externaldocument[index-]{index}

% Theorem environments.
%
\theoremstyle{plain}
\newtheorem{theorem}[subsection]{Theorem}
\newtheorem{proposition}[subsection]{Proposition}
\newtheorem{lemma}[subsection]{Lemma}

\theoremstyle{definition}
\newtheorem{definition}[subsection]{Definition}
\newtheorem{example}[subsection]{Example}
\newtheorem{exercise}[subsection]{Exercise}
\newtheorem{situation}[subsection]{Situation}

\theoremstyle{remark}
\newtheorem{remark}[subsection]{Remark}
\newtheorem{remarks}[subsection]{Remarks}

\numberwithin{equation}{subsection}

% Macros
%
\def\lim{\mathop{\mathrm{lim}}\nolimits}
\def\colim{\mathop{\mathrm{colim}}\nolimits}
\def\Spec{\mathop{\mathrm{Spec}}}
\def\Hom{\mathop{\mathrm{Hom}}\nolimits}
\def\Ext{\mathop{\mathrm{Ext}}\nolimits}
\def\SheafHom{\mathop{\mathcal{H}\!\mathit{om}}\nolimits}
\def\SheafExt{\mathop{\mathcal{E}\!\mathit{xt}}\nolimits}
\def\Sch{\mathit{Sch}}
\def\Mor{\mathop{\mathrm{Mor}}\nolimits}
\def\Ob{\mathop{\mathrm{Ob}}\nolimits}
\def\Sh{\mathop{\mathit{Sh}}\nolimits}
\def\NL{\mathop{N\!L}\nolimits}
\def\CH{\mathop{\mathrm{CH}}\nolimits}
\def\proetale{{pro\text{-}\acute{e}tale}}
\def\etale{{\acute{e}tale}}
\def\QCoh{\mathit{QCoh}}
\def\Ker{\mathop{\mathrm{Ker}}}
\def\Im{\mathop{\mathrm{Im}}}
\def\Coker{\mathop{\mathrm{Coker}}}
\def\Coim{\mathop{\mathrm{Coim}}}

% Boxtimes
%
\DeclareMathSymbol{\boxtimes}{\mathbin}{AMSa}{"02}

%
% Macros for moduli stacks/spaces
%
\def\QCohstack{\mathcal{QC}\!\mathit{oh}}
\def\Cohstack{\mathcal{C}\!\mathit{oh}}
\def\Spacesstack{\mathcal{S}\!\mathit{paces}}
\def\Quotfunctor{\mathrm{Quot}}
\def\Hilbfunctor{\mathrm{Hilb}}
\def\Curvesstack{\mathcal{C}\!\mathit{urves}}
\def\Polarizedstack{\mathcal{P}\!\mathit{olarized}}
\def\Complexesstack{\mathcal{C}\!\mathit{omplexes}}
% \Pic is the operator that assigns to X its picard group, usage \Pic(X)
% \Picardstack_{X/B} denotes the Picard stack of X over B
% \Picardfunctor_{X/B} denotes the Picard functor of X over B
\def\Pic{\mathop{\mathrm{Pic}}\nolimits}
\def\Picardstack{\mathcal{P}\!\mathit{ic}}
\def\Picardfunctor{\mathrm{Pic}}
\def\Deformationcategory{\mathcal{D}\!\mathit{ef}}
\usepackage{cleveref}
\crefname{originalreference}{original reference}{original references}
\Crefname{originalreference}{Original reference}{Original references}
\makeatletter
\newlabel{algebraic-lemma-base-change-representable-transformations-property}{{045C}{source}{}{sourcepointer-1}{}}
\newlabel{algebraic-lemma-base-change-representable-transformations-property@cref}{{[originalreference][0][]045C}{[1][0][]source}{}{}{}}
\makeatother
\crefname{originalreference}{original reference}{original references}
\Crefname{originalreference}{Original reference}{Original references}
\makeatletter
\newlabel{algebraic-lemma-composition-representable-transformations-property}{{045B}{source}{}{sourcepointer-2}{}}
\newlabel{algebraic-lemma-composition-representable-transformations-property@cref}{{[originalreference][0][]045B}{[1][0][]source}{}{}{}}
\makeatother
\crefname{originalreference}{original reference}{original references}
\Crefname{originalreference}{Original reference}{Original references}
\makeatletter
\newlabel{algebraic-lemma-equivalent}{{03YQ}{source}{}{sourcepointer-3}{}}
\newlabel{algebraic-lemma-equivalent@cref}{{[originalreference][0][]03YQ}{[1][0][]source}{}{}{}}
\makeatother
\crefname{originalreference}{original reference}{original references}
\Crefname{originalreference}{Original reference}{Original references}
\makeatletter
\newlabel{algebraic-lemma-map-space-into-stack}{{04T4}{source}{}{sourcepointer-4}{}}
\newlabel{algebraic-lemma-map-space-into-stack@cref}{{[originalreference][0][]04T4}{[1][0][]source}{}{}{}}
\makeatother
\crefname{originalreference}{original reference}{original references}
\Crefname{originalreference}{Original reference}{Original references}
\makeatletter
\newlabel{algebraic-remark-flat-fp-presentation}{{04WY}{source}{}{sourcepointer-5}{}}
\newlabel{algebraic-remark-flat-fp-presentation@cref}{{[originalreference][0][]04WY}{[1][0][]source}{}{}{}}
\makeatother
\crefname{originalreference}{original reference}{original references}
\Crefname{originalreference}{Original reference}{Original references}
\makeatletter
\newlabel{algebraic-section-2-yoneda}{{04SS}{source}{}{sourcepointer-6}{}}
\newlabel{algebraic-section-2-yoneda@cref}{{[originalreference][0][]04SS}{[1][0][]source}{}{}{}}
\makeatother
\crefname{originalreference}{original reference}{original references}
\Crefname{originalreference}{Original reference}{Original references}
\makeatletter
\newlabel{spaces-groupoids-lemma-quotient-stack-morphisms}{{044V}{source}{}{sourcepointer-7}{}}
\newlabel{spaces-groupoids-lemma-quotient-stack-morphisms@cref}{{[originalreference][0][]044V}{[1][0][]source}{}{}{}}
\makeatother
\crefname{originalreference}{original reference}{original references}
\Crefname{originalreference}{Original reference}{Original references}
\makeatletter
\newlabel{spaces-groupoids-lemma-quotient-stack-objects}{{044X}{source}{}{sourcepointer-8}{}}
\newlabel{spaces-groupoids-lemma-quotient-stack-objects@cref}{{[originalreference][0][]044X}{[1][0][]source}{}{}{}}
\makeatother
\crefname{originalreference}{original reference}{original references}
\Crefname{originalreference}{Original reference}{Original references}
\makeatletter
\newlabel{spaces-groupoids-lemma-quotient-stack-restrict-equivalence}{{046T}{source}{}{sourcepointer-9}{}}
\newlabel{spaces-groupoids-lemma-quotient-stack-restrict-equivalence@cref}{{[originalreference][0][]046T}{[1][0][]source}{}{}{}}
\makeatother
\crefname{originalreference}{original reference}{original references}
\Crefname{originalreference}{Original reference}{Original references}
\makeatletter
\newlabel{stacks-lemma-characterize-essentially-surjective-when-ff}{{046N}{source}{}{sourcepointer-10}{}}
\newlabel{stacks-lemma-characterize-essentially-surjective-when-ff@cref}{{[originalreference][0][]046N}{[1][0][]source}{}{}{}}
\makeatother
\crefname{originaltheorem}{original theorem}{original theorems}
\Crefname{originaltheorem}{Original theorem}{Original theorems}
\makeatletter
\newlabel{theorem-flat-groupoid-gives-algebraic-stack}{{06FI}{source}{}{sourcepointer-11}{}}
\newlabel{theorem-flat-groupoid-gives-algebraic-stack@cref}{{[originaltheorem][0][]06FI}{[1][0][]source}{}{}{}}
\makeatother
\title{Criterion for a quotient stack to be algebraic}
\author{The Stacks Project Authors}
\date{stacks-project:criteria.tex:local-snapshot-20261012; GNU FDL 1.2 or later}
\begin{document}
\maketitle
\noindent\textit{Editorial note.} No omission markers in-span; proof complete as written.\par\medskip
\noindent\textit{Allowed standard prerequisites (editorial).}\begin{itemize}
\item Restriction equivalence. spaces-groupoids-lemma-quotient-stack-restrict-equivalence (cross-chapter, named). [U'/R'] -> [U/R] is an equivalence under (1).
\item Flat groupoid theorem. theorem-flat-groupoid-gives-algebraic-stack (same chapter, named). Algebraicity of [U'/R'].
\item Quotient stack objects and morphisms. spaces-groupoids-lemma-quotient-stack-objects, spaces-groupoids-lemma-quotient-stack-morphisms (cross-chapter, named). Explicit description of the quotient stack over W.
\item Essential surjectivity criterion. stacks-lemma-characterize-essentially-surjective-when-ff (cross-chapter, named). Fully faithful plus local essential surjectivity.
\end{itemize}\medskip
\ifdefined\hypertarget\hypertarget{sourcepointer-1}{}\fi
\noindent\textit{Source reference.} Original reference 045C: focused-seeds/geometry/S9-crit-quotient-algebraic/ref-tags.txt, source lines 5--5. Navigation pointer only: the counted proof cites this statement by \textbackslash{}ref as a standing input; the deduction the counted proof has to make is carried out inside the counted span, and no part of that deduction is delegated to this pointer.\par\medskip
\ifdefined\hypertarget\hypertarget{sourcepointer-2}{}\fi
\noindent\textit{Source reference.} Original reference 045B: focused-seeds/geometry/S9-crit-quotient-algebraic/ref-tags.txt, source lines 6--6. Navigation pointer only: the counted proof cites this statement by \textbackslash{}ref as a standing input; the deduction the counted proof has to make is carried out inside the counted span, and no part of that deduction is delegated to this pointer.\par\medskip
\ifdefined\hypertarget\hypertarget{sourcepointer-3}{}\fi
\noindent\textit{Source reference.} Original reference 03YQ: focused-seeds/geometry/S9-crit-quotient-algebraic/ref-tags.txt, source lines 7--7. Navigation pointer only: the counted proof cites this statement by \textbackslash{}ref as a standing input; the deduction the counted proof has to make is carried out inside the counted span, and no part of that deduction is delegated to this pointer.\par\medskip
\ifdefined\hypertarget\hypertarget{sourcepointer-4}{}\fi
\noindent\textit{Source reference.} Original reference 04T4: focused-seeds/geometry/S9-crit-quotient-algebraic/ref-tags.txt, source lines 8--8. Navigation pointer only: the counted proof cites this statement by \textbackslash{}ref as a standing input; the deduction the counted proof has to make is carried out inside the counted span, and no part of that deduction is delegated to this pointer.\par\medskip
\ifdefined\hypertarget\hypertarget{sourcepointer-5}{}\fi
\noindent\textit{Source reference.} Original reference 04WY: focused-seeds/geometry/S9-crit-quotient-algebraic/ref-tags.txt, source lines 9--9. Navigation pointer only: the counted proof cites this statement by \textbackslash{}ref as a standing input; the deduction the counted proof has to make is carried out inside the counted span, and no part of that deduction is delegated to this pointer.\par\medskip
\ifdefined\hypertarget\hypertarget{sourcepointer-6}{}\fi
\noindent\textit{Source reference.} Original reference 04SS: focused-seeds/geometry/S9-crit-quotient-algebraic/ref-tags.txt, source lines 10--10. Navigation pointer only: the counted proof cites this statement by \textbackslash{}ref as a standing input; the deduction the counted proof has to make is carried out inside the counted span, and no part of that deduction is delegated to this pointer.\par\medskip
\ifdefined\hypertarget\hypertarget{sourcepointer-7}{}\fi
\noindent\textit{Source reference.} Original reference 044V: focused-seeds/geometry/S9-crit-quotient-algebraic/ref-tags.txt, source lines 11--11. Navigation pointer only: the counted proof cites this statement by \textbackslash{}ref as a standing input; the deduction the counted proof has to make is carried out inside the counted span, and no part of that deduction is delegated to this pointer.\par\medskip
\ifdefined\hypertarget\hypertarget{sourcepointer-8}{}\fi
\noindent\textit{Source reference.} Original reference 044X: focused-seeds/geometry/S9-crit-quotient-algebraic/ref-tags.txt, source lines 12--12. Navigation pointer only: the counted proof cites this statement by \textbackslash{}ref as a standing input; the deduction the counted proof has to make is carried out inside the counted span, and no part of that deduction is delegated to this pointer.\par\medskip
\ifdefined\hypertarget\hypertarget{sourcepointer-9}{}\fi
\noindent\textit{Source reference.} Original reference 046T: focused-seeds/geometry/S9-crit-quotient-algebraic/ref-tags.txt, source lines 13--13. Navigation pointer only: the counted proof cites this statement by \textbackslash{}ref as a standing input; the deduction the counted proof has to make is carried out inside the counted span, and no part of that deduction is delegated to this pointer.\par\medskip
\ifdefined\hypertarget\hypertarget{sourcepointer-10}{}\fi
\noindent\textit{Source reference.} Original reference 046N: focused-seeds/geometry/S9-crit-quotient-algebraic/ref-tags.txt, source lines 14--14. Navigation pointer only: the counted proof cites this statement by \textbackslash{}ref as a standing input; the deduction the counted proof has to make is carried out inside the counted span, and no part of that deduction is delegated to this pointer.\par\medskip
\ifdefined\hypertarget\hypertarget{sourcepointer-11}{}\fi
\noindent\textit{Source reference.} Original theorem 06FI: sources/stacks-criteria/pinned-original-source.tex, source lines 3245--3245. Navigation pointer only: the counted proof cites this statement by \textbackslash{}ref as a standing input; the deduction the counted proof has to make is carried out inside the counted span, and no part of that deduction is delegated to this pointer.\par\medskip
\begin{lemma}
\label{lemma-quotient-algebraic}
Let $S$ be a scheme and let $B$ be an algebraic space over $S$.
Let $(U, R, s, t, c)$ be a groupoid in algebraic spaces over $B$.
The quotient stack $[U/R]$ is an algebraic stack if and only if
there exists a morphism of algebraic spaces $g : U' \to U$ such that
\begin{enumerate}
\item the composition
$U' \times_{g, U, t} R \to R \xrightarrow{s} U$ is a surjection of
sheaves, and
\item the morphisms $s', t' : R' \to U'$ are flat and locally of finite
presentation where $(U', R', s', t', c')$ is the restriction of
$(U, R, s, t, c)$ via $g$.
\end{enumerate}
\end{lemma}

\begin{proof}
First, assume that $g : U' \to U$ satisfies (1) and (2). Property (1)
implies that $[U'/R'] \to [U/R]$ is an equivalence, see
Groupoids in Spaces,
Lemma \ref{spaces-groupoids-lemma-quotient-stack-restrict-equivalence}.
By
Theorem \ref{theorem-flat-groupoid-gives-algebraic-stack}
the quotient stack $[U'/R']$ is an algebraic stack. Hence
$[U/R]$ is an algebraic stack too, see
Algebraic Stacks, Lemma \ref{algebraic-lemma-equivalent}.

\medskip\noindent
Conversely, assume that $[U/R]$ is an algebraic stack. We may choose a
scheme $W$ and a surjective smooth $1$-morphism
$$
f : (\Sch/W)_{fppf} \longrightarrow [U/R].
$$
By the $2$-Yoneda lemma
(Algebraic Stacks, Section \ref{algebraic-section-2-yoneda})
this corresponds to an object $\xi$ of $[U/R]$ over $W$.
By the description of $[U/R]$ in
Groupoids in Spaces, Lemma \ref{spaces-groupoids-lemma-quotient-stack-objects}
we can find a surjective, flat, locally finitely presented morphism
$b : U' \to W$ of schemes such that $\xi' = b^*\xi$ corresponds to a morphism
$g : U' \to U$. Note that the $1$-morphism
$$
f' : (\Sch/U')_{fppf} \longrightarrow [U/R].
$$
corresponding to $\xi'$ is surjective, flat, and locally of finite
presentation, see
Algebraic Stacks, Lemma
\ref{algebraic-lemma-composition-representable-transformations-property}.
Hence
$(\Sch/U')_{fppf} \times_{[U/R]} (\Sch/U')_{fppf}$
which is represented by the algebraic space
$$
\mathit{Isom}_{[U/R]}(\text{pr}_0^*\xi', \text{pr}_1^*\xi') =
(U' \times_S U')
\times_{(g \circ \text{pr}_0, g \circ \text{pr}_1), U \times_S U} R = R'
$$
(see
Groupoids in Spaces, Lemma
\ref{spaces-groupoids-lemma-quotient-stack-morphisms}
for the first equality; the second is the definition of restriction)
is flat and locally of finite presentation over $U'$ via both $s'$ and $t'$
(by base change, see
Algebraic Stacks, Lemma
\ref{algebraic-lemma-base-change-representable-transformations-property}).
By this description of $R'$ and by
Algebraic Stacks, Lemma \ref{algebraic-lemma-map-space-into-stack}
we obtain a canonical fully faithful $1$-morphism $[U'/R'] \to [U/R]$.
This $1$-morphism is essentially surjective because $f'$ is flat,
locally of finite presentation, and surjective (see
Stacks, Lemma \ref{stacks-lemma-characterize-essentially-surjective-when-ff});
another way to prove this is to use
Algebraic Stacks, Remark
\ref{algebraic-remark-flat-fp-presentation}.
Finally, we can use
Groupoids in Spaces, Lemma
\ref{spaces-groupoids-lemma-quotient-stack-restrict-equivalence}
to conclude that the composition
$U' \times_{g, U, t} R \to R \xrightarrow{s} U$ is a surjection of sheaves.
\end{proof}

\end{document}
